\documentclass[11pt]{article}
\usepackage[margin=1in]{geometry}
\usepackage{amsmath,amssymb,amsthm}
\usepackage[colorlinks=true,linkcolor=blue,urlcolor=blue,citecolor=blue]{hyperref}
\newtheorem{theorem}{Theorem}
\newtheorem{lemma}[theorem]{Lemma}
\newtheorem{corollary}[theorem]{Corollary}
\newtheorem{proposition}[theorem]{Proposition}
\theoremstyle{remark}
\newtheorem*{remark}{Remark}
\newcommand{\C}{\mathbb C}
\newcommand{\N}{\mathbb N}
\newcommand{\E}{\mathbb E}
\newcommand{\Prob}{\mathbb P}
\newcommand{\dd}{\,d}
\newcommand{\avint}{\frac{1}{2\pi}\int_0^{2\pi}}
\DeclareMathOperator{\supp}{supp}
\title{Cofinite derivative zeros at the order-one threshold}
\author{Jensen Kohlmeyer \and Liam Kruer}
\date{October 4, 2026}
\begin{document}
\maketitle

\begin{abstract}
We give a complete probabilistic construction of a transcendental entire
function for which every nonempty open set contains a zero of every
sufficiently high derivative. The construction has
$\log M(r)=O(r\log r\log\log r)$, hence order exactly one and infinite
exponential type. No transcendental entire function of finite exponential
type can have this zero property. We also establish a local limiting
zero-counting measure for the constructed derivatives.
The original existence assertion already has prior solution claims,
including Hou's 2026 preprint. This draft does not claim first discovery
of that assertion. The order-one construction and growth obstruction are
the proposed strengthening; their novelty remains subject to literature
review. All mathematical arguments are supplied below.
\end{abstract}

\section{Problem, scope, and attribution}

The prize catalogue lists JSP-000752 as Open on the source snapshot read
on October 4, 2026 \cite{JSP}. It asks for an entire function with dense
derivative zero sets along every infinite subsequence of derivative
orders. The natural, nontrivial interpretation requires a
\emph{transcendental} entire function: a nonzero polynomial would satisfy
the literal version because all sufficiently high derivatives vanish
identically. This is Erd\H{o}s Problem 906 in Bloom's database
\cite{Bloom}.

The prize rules permit a complete mathematical solution to be supplied
before Lean formalization \cite{Rules}. Its catalogue flag
``Eligible to claim'' has a different meaning: an award claim requires
both an accepted mathematical solution and verified formalization. This manuscript supplies mathematical proof evidence for review;
no award entitlement or maintainer acceptance is asserted.

The bounded random-series and zero-free-disk framework has a clear
precedent in Hou \cite{Hou}, whose stated construction has order two.
Earlier probabilistic proposals are linked in the discussion of
Problem 906 \cite{Bloom}. The existence answer itself is therefore
\emph{not claimed as a new breakthrough}. The construction below replaces
the Fock weights with slowly increasing derivative weights and proves a
stronger growth conclusion. It uses an inverse-amplitude moment and
Jensen's inequality; no Laguerre asymptotics or external random-zero
concentration theorem is needed. The order-one strengthening was not located in the cited prior sources;
this is not a claim of verified first-discovery priority.

Write $B(c,R)=\{z:|z-c|<R\}$ and
$M_f(r)=\max_{|z|=r}|f(z)|$. The zero property we will prove is
\begin{equation}\tag{C}\label{eq:C}
 \forall U\subset\C\text{ nonempty and open},\quad
 \exists N_U\quad\forall n\geq N_U,\quad
 \exists z\in U:\ f^{(n)}(z)=0.
\end{equation}
All derivative orders are nonnegative integers.

\begin{lemma}[The required quantifier]
An entire function has property \eqref{eq:C} if and only if the union of
the zero sets of $f^{(n_k)}$ is dense for every strictly increasing
sequence $(n_k)$ of derivative orders.
\end{lemma}
\begin{proof}
If \eqref{eq:C} holds, every such sequence eventually exceeds $N_U$, so
the union meets each nonempty open $U$. Conversely, if \eqref{eq:C}
fails for one $U$, infinitely many derivative orders have no zero in
$U$. Enumerating those orders gives a strictly increasing sequence
whose union of zero sets misses $U$.
\end{proof}

\section{A sharp order threshold}

\begin{theorem}\label{thm:main}
There is a transcendental entire function $f$ with property
\eqref{eq:C} such that
\begin{equation}\label{eq:growth}
 \log M_f(r)=O(r\log r\log\log r)\qquad(r\longrightarrow\infty).
\end{equation}
It has order exactly one and infinite exponential type.
Moreover, no transcendental entire function satisfying
\eqref{eq:C} has finite exponential type. Thus one is the smallest
possible order, and this order is attained.
\end{theorem}

Finite exponential type here means that
$|g(z)|\leq C e^{T|z|}$ for some finite $C,T>0$ and all $z$.
The usual order is
$\rho(g)=\limsup_{r\to\infty}\log\log M_g(r)/\log r$.

\begin{proposition}[Finite exponential type is impossible]
\label{prop:obstruction}
A transcendental entire function of finite exponential type cannot
satisfy \eqref{eq:C}.
\end{proposition}
\begin{proof}
Suppose $|g(z)|\leq C e^{T|z|}$ with $T>0$. Fix $R>0$, and let
\[
 m_n(R)=\max_{|z|\leq R}|g^{(n)}(z)|.
\]
If \eqref{eq:C} holds, for all $n\geq N_R$ choose
$a_n\in B(0,R)$ with $g^{(n)}(a_n)=0$.
The line segment from $a_n$ to any $z$ in the closed disk lies in
that disk. The fundamental theorem of calculus along this segment gives
\[
 |g^{(n)}(z)|\leq |z-a_n|m_{n+1}(R)\leq2R m_{n+1}(R).
\]
Consequently
\begin{equation}\label{eq:lowerder}
 m_n(R)\geq m_{N_R}(R)(2R)^{-(n-N_R)}\qquad(n\geq N_R).
\end{equation}
The factor $m_{N_R}(R)$ is strictly positive, since a derivative
of a transcendental entire function is not identically zero.

Cauchy's estimate on a circle of radius $t$ about $z$, $|z|\leq R$,
gives
\[
 m_n(R)\leq C e^{TR} n!\frac{e^{Tt}}{t^n}.
\]
Take $t=n/T$ and use $n!\leq n^n$. For $n\geq1$ this yields
\begin{equation}\label{eq:upperder}
 m_n(R)\leq C e^{TR}(eT)^n.
\end{equation}
Choosing $0<R<(2eT)^{-1}$ contradicts
\eqref{eq:lowerder}--\eqref{eq:upperder}, because
$(2R)^{-1}>eT$.
\end{proof}

\section{The weighted random entire function}

Set $A=e^e$ and, for $t\geq0$, define
\begin{equation}\label{eq:w}
 w(t)=\log(t+A)\log\log(t+A),\qquad W_n=w(n).
\end{equation}
Define $b_0=1$ and $b_n=\prod_{j=1}^n W_j$ for $n\geq1$.
Let $(\xi_j)_{j\geq0}$ be independent random variables, each uniformly
distributed with respect to area on the closed unit disk.
Our entire function is
\begin{equation}\label{eq:f}
 f(z)=\sum_{j=0}^\infty\xi_j\frac{b_j}{j!}z^j.
\end{equation}
These are genuine random complex coefficients; the probabilistic
argument extracts one deterministic sample at the end.

We will use only the following elementary facts about $w$:
\begin{equation}\label{eq:wfacts}
 W_n\geq1,\quad w\text{ is increasing},\quad
 \frac{W_{2n}}{W_n}\longrightarrow1,\quad
 \frac{W_n}{\log n}\longrightarrow\infty,\quad
 \frac{W_n}{\sqrt n}\longrightarrow0.
\end{equation}
Indeed, for $t\geq0$,
$w'(t)=(\log\log(t+A)+1)/(t+A)>0$, and the remaining
limits follow from $w(t)\sim\log t\log\log t$.

\begin{lemma}[Normal convergence and deterministic growth]
\label{lem:entire}
Every sample in \eqref{eq:f} is entire. For all sufficiently large $r$,
uniformly over all samples,
\begin{equation}\label{eq:explicitgrowth}
 M_f(r)\leq2\exp\!\bigl(rw(\lceil r^2\rceil)\bigr).
\end{equation}
Almost surely all Taylor coefficients are nonzero, so $f$ is
transcendental almost surely.
\end{lemma}
\begin{proof}
Monotonicity gives $b_j\leq W_j^j$. The elementary factorial bound
$j!\geq(j/e)^j$ gives
\[
 \left(\frac{b_j}{j!}\right)^{1/j}\leq\frac{eW_j}{j}
 \longrightarrow0.
\]
Thus the series converges normally on compact sets, uniformly in the
bounded coefficients. The same is true after any fixed number of
termwise differentiations.

For the growth bound put $J=\lceil r^2\rceil$.
For $0\leq j\leq J$ we have $b_j\leq W_J^j$, so the sum of
these absolute terms is at most $\exp(rW_J)$.
For $j>J$, $r\leq\sqrt j$, and therefore
\[
 \frac{b_jr^j}{j!}
 \leq\left(\frac{erW_j}{j}\right)^j
 \leq\left(\frac{eW_j}{\sqrt j}\right)^j.
\]
By \eqref{eq:wfacts}, for all sufficiently large $r$ the last
quantity is at most $2^{-j}$ for every $j>J$.
The tail is at most $2^{-J}\leq1$, which is at most the preceding
exponential bound. This proves \eqref{eq:explicitgrowth}.
Finally each coefficient law is nonatomic, so
$\Prob(\xi_j=0)=0$ for each $j$. A countable union of null events
still has probability zero.
\end{proof}

Normalize the derivatives by
\begin{equation}\label{eq:G}
 G_n(z)=\frac{f^{(n)}(z)}{b_n}
 =\sum_{k=0}^\infty\xi_{n+k}B_{n,k}\frac{z^k}{k!},\qquad
 B_{n,k}=\frac{b_{n+k}}{b_n}
 =\prod_{j=n+1}^{n+k}W_j.
\end{equation}
The empty product is one. Normalization does not change any zero set.

\begin{lemma}[A uniform upper bound for a derivative]
\label{lem:upper}
For every fixed $H>0$, for all sufficiently large $n$ and all
$0\leq r\leq H$,
\begin{equation}\label{eq:Sn}
 S_n(r):=\sum_{k=0}^\infty B_{n,k}\frac{r^k}{k!}
 \leq e^{rW_{2n}}+2^{-n}\leq2e^{rW_{2n}}.
\end{equation}
In particular $|G_n(z)|\leq2e^{|z|W_{2n}}$ for $|z|\leq H$.
\end{lemma}
\begin{proof}
For $k\leq n$, all factors in $B_{n,k}$ are at most $W_{2n}$.
These terms sum to at most $e^{rW_{2n}}$.
For $k>n$, we have $n+k\leq2k$, hence
\[
 B_{n,k}\frac{r^k}{k!}
 \leq\left(\frac{eH w(2k)}{k}\right)^k.
\]
Since $w(2k)/k\to0$, for all sufficiently large $n$ the last
quantity is at most $2^{-k}$ simultaneously for every $k>n$.
The sum of these terms is $2^{-n}$. The last assertion uses
$|\xi_j|\leq1$.
\end{proof}

\section{One-coordinate anti-concentration}

\begin{lemma}[An inverse-amplitude moment]\label{lem:inverse}
If $\xi$ is uniform on the unit disk, $a\neq0$ and $v\in\C$, then
\begin{equation}\label{eq:smallball}
 \Prob(|a\xi+v|<s)\leq\min\{1,s^2/|a|^2\}\quad(s>0),
 \qquad \E\frac{1}{|a\xi+v|}\leq\frac{2}{|a|}.
\end{equation}
\end{lemma}
\begin{proof}
The first inequality follows by comparing areas: the preimage of the
disk of radius $s$ has area at most $\pi s^2/|a|^2$ in the unit disk.
The layer-cake formula and this inequality give
\[
 \E\frac1{|a\xi+v|}
 \leq\int_0^{1/|a|}1\dd t+
 \int_{1/|a|}^\infty\frac{1}{|a|^2t^2}\dd t
 =\frac{2}{|a|}.
\]
\end{proof}

\begin{lemma}[A near-peak coefficient]\label{lem:peak}
For $t\geq1$ and $m=\lfloor t\rfloor$,
\begin{equation}\label{eq:peakfactorial}
 \frac{t^m}{m!}\geq\frac{e^t}{e^2t}.
\end{equation}
Consequently, if $|z|W_n\geq1$, the term with
$m=\lfloor |z|W_n\rfloor$ in \eqref{eq:G} has coefficient
$a=B_{n,m}z^m/m!$ satisfying
\[
 |a|\geq \frac{e^{|z|W_n}}{e^2|z|W_n}.
\]
\end{lemma}
\begin{proof}
For $m\geq1$, comparison with an integral gives
\[
 \log(m!)\leq m\log m-m+1+\log m.
\]
Thus
\[
 \log\frac{t^m}{m!}
 \geq m\log(t/m)+m-1-\log m
 \geq t-2-\log t,
\]
because $m\leq t$, $m\geq t-1$, and $\log m\leq\log t$.
Monotonicity of $w$ gives $B_{n,m}\geq W_n^m$, proving
the second assertion.
\end{proof}

\begin{lemma}[Uniform inverse moments on a compact set]
\label{lem:pointmoment}
For any $H>0$, put $H_+=\max\{1,H\}$. For every $n\geq1$ and
$|z|\leq H$,
\begin{equation}\label{eq:pointmoment}
 \E\left(\frac{e^{W_n|z|}}{|G_n(z)|}\right)
 \leq 2e^2H_+W_n.
\end{equation}
The value at $G_n(z)=0$ is interpreted as $+\infty$.
\end{lemma}
\begin{proof}
If $|z|W_n\geq1$, choose $m$ as in Lemma~\ref{lem:peak} and
condition on all coefficients except $\xi_{n+m}$. The normally
convergent series becomes $a\xi_{n+m}+v$, with $v$ fixed under this
conditioning. Independence, Lemma~\ref{lem:inverse}, and
Lemma~\ref{lem:peak} give
\[
 \E\left(\left.\frac{e^{W_n|z|}}{|G_n(z)|}\right|
                  \text{all other coefficients}\right)
 \leq2e^2|z|W_n\leq2e^2H_+W_n.
\]
Integration removes the conditioning.
If $|z|W_n<1$, condition instead on $\xi_n$, whose coefficient
in $G_n(z)$ is one. The left side is then at most $2e$,
which is no larger than the right side since $H_+,W_n\geq1$.
This includes $z=0$.
\end{proof}

\section{Summable probabilities of zero-free disks}

\begin{lemma}[Strict circular gain]\label{lem:gain}
For $c\in\C$ and $\eta>0$,
\[
 \gamma(c,\eta)=\avint|c+\eta e^{i\theta}|\dd\theta-|c|>0.
\]
\end{lemma}
\begin{proof}
Pair the angles $\theta$ and $\theta+\pi$. The triangle inequality
gives $|c+\eta e^{i\theta}|+|c-\eta e^{i\theta}|\geq2|c|$.
For $c\neq0$ the inequality is strict on an interval of angles near
a direction perpendicular to $c$. For $c=0$ the assertion is
immediate. Integrating proves strict positivity.
\end{proof}

Fix $c\neq0$ and $0<R<|c|$, put $\eta=R/2$ and
$H=|c|+\eta$, and let $\mathcal H_n(c,R)$ be the event that
$G_n$ has no zero in $B(c,R)$. Define the nonnegative random variable
\begin{equation}\label{eq:Y}
 Y_n=\avint
 \frac{\exp(W_n|c+\eta e^{i\theta}|)}
 {|G_n(c+\eta e^{i\theta})|}\dd\theta.
\end{equation}
Tonelli's theorem and \eqref{eq:pointmoment} imply
\begin{equation}\label{eq:Ymean}
 \E Y_n\leq2e^2H_+W_n.
\end{equation}

On $\mathcal H_n(c,R)$, $\log|G_n|$ is harmonic throughout
$B(c,R)$, and therefore
\begin{equation}\label{eq:harmonicmean}
 \avint\log|G_n(c+\eta e^{i\theta})|\dd\theta
 =\log|G_n(c)|.
\end{equation}
Jensen's inequality for the exponential, applied to \eqref{eq:Y},
and \eqref{eq:harmonicmean} give
\[
 \log Y_n\geq W_n(|c|+\gamma(c,\eta))-\log|G_n(c)|.
\]
For all sufficiently large $n$, Lemma~\ref{lem:upper} bounds the
last term. Thus, on $\mathcal H_n(c,R)$,
\begin{equation}\label{eq:holeimplies}
 Y_n\geq\frac12
 \exp\!\left(\gamma(c,\eta)W_n-|c|(W_{2n}-W_n)\right).
\end{equation}
Markov's inequality, \eqref{eq:Ymean}, and
\eqref{eq:holeimplies} now yield
\begin{align}
 \Prob(\mathcal H_n(c,R))
 &\leq4e^2H_+W_n
  \exp\!\left(-\gamma(c,\eta)W_n+|c|(W_{2n}-W_n)\right)
   \label{eq:hole1}\\
 &\leq4e^2H_+W_n\exp\!\left(-\tfrac12\gamma(c,\eta)W_n\right)
   \label{eq:hole2}
\end{align}
for all sufficiently large $n$. The second inequality uses
$W_{2n}/W_n\to1$ and $\gamma(c,\eta)>0$.

For every $a>0$, \eqref{eq:wfacts} implies
\begin{equation}\label{eq:summable}
 \sum_{n=1}^\infty W_n e^{-aW_n}<\infty.
\end{equation}
For example, eventually $W_n\leq\sqrt n$ and
$aW_n\geq4\log n$, so the summand is at most $n^{-7/2}$.
Consequently the probabilities in \eqref{eq:hole2} are summable.
The first Borel--Cantelli lemma proves that, almost surely,
only finitely many $G_n$ are zero-free in this fixed disk.
No independence between different derivative orders is used.

For completeness the events above are measurable. Normal convergence
makes $(\omega,z)\mapsto G_n(\omega,z)$ jointly continuous when the
sample space is the countable product of closed unit disks with its
product topology. Zero-freeness on an open disk equals positivity of
the minimum modulus on each of a countable increasing sequence of
closed subdisks. Each minimum is a measurable function; it can also
be computed as an infimum over a countable dense set. The nonnegative
integrand in \eqref{eq:Y} is measurable, so Tonelli applies without
any integrability assumption in advance.

\begin{proof}[Completion of Theorem~\ref{thm:main}]
Take the countable family
\[
 \mathcal B=\{B(c,R):c\in\mathbb Q+i\mathbb Q,\quad
                    R\in\mathbb Q_{>0},\quad R<|c|\}.
\]
Every nonempty open subset of $\C$ contains a member of $\mathcal B$:
choose a nonzero point in it, then a nearby rational complex center
and a sufficiently small rational radius.
The preceding Borel--Cantelli conclusion holds simultaneously for
all disks in $\mathcal B$ outside a countable union of null events.
Intersect with the probability-one event on which all coefficients
are nonzero. Choose a sample in this intersection.
The resulting deterministic transcendental entire $f$ has
\eqref{eq:C}; the quantifier lemma gives the original subsequence
formulation as well.

Lemma~\ref{lem:entire} implies \eqref{eq:growth}, since
$w(\lceil r^2\rceil)=O(\log r\log\log r)$.
Thus $\rho(f)\leq1$.
Every entire function of order less than one has finite exponential
type: choose $\rho(f)<q<1$, then eventually
$\log M_f(r)\leq r^q\leq r$, and adjust a constant on the remaining
compact disk. Proposition~\ref{prop:obstruction} therefore forces
$\rho(f)\geq1$. It also rules out finite exponential type for our
chosen function. This proves all assertions of the theorem.
\end{proof}

\section{A further check: limiting zero density}

The same construction provides a quantitative version. This section
is not needed for the existence proof above.
Let $\nu_n$ be the locally finite measure that counts the zeros of
$f^{(n)}$, with multiplicities.

\begin{theorem}\label{thm:density}
For almost every sample in \eqref{eq:f},
\[
 \frac{\nu_n}{W_n}\ \longrightarrow\
 \frac{1}{2\pi|z|}\dd A(z)
\]
locally weakly as $n\to\infty$. The limiting measure has no atom at
zero; its displayed density is locally integrable.
In particular, for every fixed $r>0$,
\[
 \nu_n(B(0,r))\sim rW_n.
\]
\end{theorem}
\begin{proof}
For a fixed $H>0$, write
\[
 V_n(z)=\bigl(W_n|z|-\log|G_n(z)|\bigr)_+,
 \qquad Q_n=\frac1{\pi H^2}\int_{|z|<H}V_n(z)\dd A(z).
\]
The value $+\infty$ at a zero causes no difficulty. Nonzero analytic
functions have locally integrable logarithms. Moreover
$e^{x_+}\leq1+e^x$ and Lemma~\ref{lem:pointmoment} give
\[
 \E e^{V_n(z)}\leq1+2e^2H_+W_n\qquad(|z|\leq H).
\]
Jensen's inequality for normalized area measure and Tonelli's theorem
imply
\[
 \E e^{Q_n}\leq1+2e^2H_+W_n.
\]
Hence, for every $\varepsilon>0$,
\[
 \Prob(Q_n>\varepsilon W_n)
 \leq(1+2e^2H_+W_n)e^{-\varepsilon W_n}.
\]
These probabilities are summable by \eqref{eq:summable}.
Apply Borel--Cantelli for integer $H\geq1$ and positive rational
$\varepsilon$ to obtain, almost surely, $Q_n/W_n\to0$ on every
compact disk.

Set $u_n(z)=W_n^{-1}\log|G_n(z)|$. Lemma~\ref{lem:upper} gives
\[
 (u_n(z)-|z|)_+\leq
 H\left(\frac{W_{2n}}{W_n}-1\right)+\frac{\log2}{W_n}
 \longrightarrow0
\]
uniformly for $|z|\leq H$. The opposite part is exactly
$(|z|-u_n(z))_+=V_n(z)/W_n$. Thus
\begin{equation}\label{eq:L1}
 u_n\longrightarrow |z|\quad\hbox{in }L^1_{\rm loc}(\C)
 \quad\hbox{almost surely}.
\end{equation}
For a nonzero analytic function $h$, the elementary distributional
identity
\[
 \Delta\log|h|=2\pi\sum_{h(a)=0}\operatorname{mult}_a(h)\delta_a
\]
follows from local factorization at each zero and
$\Delta\log|z-a|=2\pi\delta_a$; the nonvanishing local factor
has harmonic log modulus. Therefore
$\nu_n/W_n=(2\pi)^{-1}\Delta u_n$.
Taking distributional Laplacians in \eqref{eq:L1} gives
\[
 \frac{\nu_n}{W_n}\longrightarrow\frac1{2\pi}\Delta|z|
 =\frac1{2\pi|z|}\dd A(z).
\]
For the last equality, away from zero the radial Laplacian is
$1/|z|$. Its flux through a circle of radius $s$ is $2\pi s$,
which tends to zero, so there is no delta mass at the origin.

Distributional convergence of these positive measures implies local
weak convergence: a smooth compactly supported upper cutoff bounds
their masses on any fixed compact set, and continuous test functions
can then be uniformly approximated by smooth ones. Each disk boundary
has zero limiting measure, so the masses of any fixed disk converge
to the mass prescribed by the limit. Polar integration gives
$\int_{B(0,r)}(2\pi|z|)^{-1}\dd A=r$.
\end{proof}

\section{General weights and verification notes}

The proof uses no special identity for \eqref{eq:w}. It works for any
nondecreasing sequence $W_n\geq1$ with
\[
 W_{2n}/W_n\to1,\qquad W_n/\log n\to\infty,\qquad
 W_n/\sqrt n\to0,
\]
using $b_n=\prod_{j=1}^nW_j$. The deterministic bound becomes
$\log M_f(r)\leq rW_{\lceil r^2\rceil}+\log2$ eventually,
and the limiting zero measure has normalization $W_n$.
For example, slowly varying factors growing more slowly than
$\log\log n$ can replace that factor in \eqref{eq:w}, as long as
the three displayed conditions hold. The exact least possible
growth rate inside the order-one infinite-type class is not determined
by this argument.

The main proof is analytic and has no numerical assumptions. An
accompanying Python script checks the coefficient peak inequality and
the derivative majorant at selected parameter values using high
precision arithmetic. Those checks are sanity checks, not substitutes
for the proofs and not a Lean formalization.

\section*{AI assistance and authorship disclosure}

Jensen Kohlmeyer (GitHub: \href{https://github.com/JENW1N}{JENW1N})
and Liam Kruer (GitHub: \href{https://github.com/lkruer}{lkruer})
are the designated authors of this AI-assisted research manuscript.
OpenAI Codex assisted with the mathematical construction, proof development,
coefficient checks, literature comparison, and manuscript preparation.
The public submission is made from the JENW1N account.
No independent human verification or distinct individual author roles
are asserted by this disclosure; attribution and correctness are
submitted for review.

The principal points for independent review are the conditional inverse
moment, the uniform tail estimate for all $k>n$, and the summability
for every positive disk-dependent constant. The full proof above
addresses the original infinite-subsequence quantifier. It has no
unproved mathematical premise. First-solution credit is expressly not
claimed; novelty of the sharper growth statements remains unverified.

\begin{thebibliography}{9}
\bibitem{JSP} The Justin Sun Prize, problem JSP-000752, catalogue
snapshot read October 4, 2026.
\url{https://github.com/TheJustinSunPrize/awards/blob/main/problems/catalog-0701-0800.md}
\bibitem{Rules} The Justin Sun Prize, contribution guidelines,
snapshot read October 4, 2026.
\url{https://github.com/TheJustinSunPrize/awards/blob/main/CONTRIBUTING.md}
\bibitem{Bloom} T. F. Bloom, Erd\H{o}s Problem 906 and its discussion.
\url{https://www.erdosproblems.com/906}
\url{https://www.erdosproblems.com/forum/thread/906?embed=1}
\bibitem{Hou} E. Hou, \emph{Cofinite Zeros of High Derivatives},
arXiv:2607.20816, 2026. Version 3 and the author's current repository
manuscript were inspected for this draft.
\url{https://arxiv.org/html/2607.20816v3}
\url{https://github.com/erichou1/cofinite-derivative-zeros}
\end{thebibliography}
\end{document}
